% GENERATED FILE. Edit canonical lessons, then run npm run generate:books.
% canonical-source-hash: fnv1a64:efb3b5997c85b909
% canonical-lessons: FA-C322-mahkum-shodan, FA-C322-tabrae-shodan, FA-C322-zendani, FA-C322-vasiqe, FA-C322-azad-shodan

\chapter{Verdict and release}
\label{ch:fa-verdict-and-release}
\hlchaptermodality{322}

\begin{chapteropening}
\textbf{By the end of this chapter:} \emph{I can say that someone was sentenced or cleared, held in prison, let out on bail and released.}

Complete the last lesson of chapter 322: I can say that someone was sentenced or cleared, held in prison, let out on bail and released.
\end{chapteropening}

\section[\texorpdfstring{\fa{محکوم} \fa{شدن} (mahkum shodan)}{محکوم شدن (mahkum shodan)}]{\fa{محکوم} \fa{شدن} (mahkum shodan) — to be sentenced, to be convicted}

\label{lesson:FA-C322-mahkum-shodan}



\begin{quote}
Before the new one: say the Persian for to give away, then the Persian for guilty.
\end{quote}

\subsection*{You'll want to know: \fa{محکوم} \fa{شدن}}
\textbf{\fa{محکوم} \fa{شدن}} (\emph{mahkum shodan}) — \textquotedblleft{}to be sentenced, to be convicted\textquotedblright{}.

\begin{quote}
\fa{مرد} \fa{به} \fa{دو} \fa{سال} \fa{زندان} \fa{محکوم} \fa{شد}.
\end{quote}

\begin{quote}
\emph{Mard be do sâl zendân mahkum shod.}
\end{quote}

\begin{quote}
The man was sentenced to two years in prison.
\end{quote}

The punishment takes \emph{be}. Present \emph{mahkum mishavad}.

\begin{cousinweb}
\emph{mahkum} is usually traced to Arabic \emph{maḥkūm}, \textquotedblleft{}judged, ruled on\textquotedblright{}; EtymDB lists no entry. \emph{shodan} \textquotedblleft{}to become\textquotedblright{} is inherited Persian.
\end{cousinweb}

\begin{grammarlens}[title={What you've built}]
One more word for telling, explaining and describing at B1.
\end{grammarlens}

\subsection*{Your turn}
\begin{itemize}
\raggedright
  \item \emph{Say it:} \emph{mahkum shodan}
  \item \emph{Say it:} \emph{mahkum shodan}, once more
  \item \emph{Say it:} a sentence of your own with \emph{mahkum shodan}
  \item \emph{Recall:} say the Persian for to give away, then the Persian for guilty, then say \emph{mahkum shodan} again
\end{itemize}

\begin{tcolorbox}[breakable,colback=teal!4,colframe=teal!35!black,title={Before you move on}]
What does \fa{محکوم} \fa{شدن} mean? (\textquotedblleft{}To be sentenced, to be convicted\textquotedblright{}.) Say it once more.
\end{tcolorbox}
\section[\texorpdfstring{\fa{تبرئه} \fa{شدن} (tabra'e shodan)}{تبرئه شدن (tabra'e shodan)}]{\fa{تبرئه} \fa{شدن} (tabra'e shodan) — to be acquitted, to be cleared}

\label{lesson:FA-C322-tabrae-shodan}



\begin{quote}
Before the new one: say the Persian for guilty, then the Persian for to be sentenced.
\end{quote}

\subsection*{You'll want to know: \fa{تبرئه} \fa{شدن}}
\textbf{\fa{تبرئه} \fa{شدن}} (\emph{tabra'e shodan}) — \textquotedblleft{}to be acquitted, to be cleared\textquotedblright{}.

\begin{quote}
\fa{او} \fa{بعد} \fa{از} \fa{یک} \fa{سال} \fa{تبرئه} \fa{شد}.
\end{quote}

\begin{quote}
\emph{U ba'd az yek sâl tabra'e shod.}
\end{quote}

\begin{quote}
She was cleared after a year.
\end{quote}

Written with the hamza seat \fa{ئ}. Present \emph{tabra'e mishavad}.

\begin{cousinweb}
\emph{tabra'e} is usually traced to Arabic \emph{tabri'a}, \textquotedblleft{}a clearing of guilt\textquotedblright{}; EtymDB lists no entry. \emph{shodan} \textquotedblleft{}to become\textquotedblright{} is inherited Persian.
\end{cousinweb}

\begin{grammarlens}[title={What you've built}]
One more word for telling, explaining and describing at B1.
\end{grammarlens}

\subsection*{Your turn}
\begin{itemize}
\raggedright
  \item \emph{Say it:} \emph{tabra'e shodan}
  \item \emph{Say it:} \emph{tabra'e shodan}, once more
  \item \emph{Say it:} a sentence of your own with \emph{tabra'e shodan}
  \item \emph{Recall:} say the Persian for guilty, then the Persian for to be sentenced, then say \emph{tabra'e shodan} again
\end{itemize}

\begin{tcolorbox}[breakable,colback=teal!4,colframe=teal!35!black,title={Before you move on}]
What does \fa{تبرئه} \fa{شدن} mean? (\textquotedblleft{}To be acquitted, to be cleared\textquotedblright{}.) Say it once more.
\end{tcolorbox}
\section[\texorpdfstring{\fa{زندانی} (zendâni)}{زندانی (zendâni)}]{\fa{زندانی} (zendâni) — a prisoner}

\label{lesson:FA-C322-zendani}



\begin{quote}
Before the new one: say the Persian for to be sentenced, then the Persian for to be acquitted.
\end{quote}

\subsection*{You'll want to know: \fa{زندانی}}
\textbf{\fa{زندانی}} (\emph{zendâni}) — \textquotedblleft{}a prisoner\textquotedblright{}.

\begin{quote}
\fa{زندانی} \fa{ها} \fa{هر} \fa{روز} \fa{یک} \fa{ساعت} \fa{در} \fa{حیاط} \fa{بودند}.
\end{quote}

\begin{quote}
\emph{Zendâni-hâ har ruz yek sâ'at dar hayât budand.}
\end{quote}

\begin{quote}
The prisoners spent an hour in the yard every day.
\end{quote}

\emph{zendân} \textquotedblleft{}prison\textquotedblright{} + \emph{-i}. Also an adjective, \textquotedblleft{}imprisoned\textquotedblright{}.

\begin{cousinweb}
\emph{zendâni} adds the suffix \emph{-i} to \emph{zendân} \textquotedblleft{}prison\textquotedblright{}, a native Persian word; EtymDB links \emph{zendân} to Old Armenian \emph{zndan} but gives it no parent, so its deeper origin is not settled.
\end{cousinweb}

\begin{grammarlens}[title={What you've built}]
One more word for telling, explaining and describing at B1.
\end{grammarlens}

\subsection*{Your turn}
\begin{itemize}
\raggedright
  \item \emph{Say it:} \emph{zendâni}
  \item \emph{Say it:} \emph{zendâni}, once more
  \item \emph{Say it:} a sentence of your own with \emph{zendâni}
  \item \emph{Recall:} say the Persian for to be sentenced, then the Persian for to be acquitted, then say \emph{zendâni} again
\end{itemize}

\begin{tcolorbox}[breakable,colback=teal!4,colframe=teal!35!black,title={Before you move on}]
What does \fa{زندانی} mean? (\textquotedblleft{}A prisoner\textquotedblright{}.) Say it once more.
\end{tcolorbox}
\section[\texorpdfstring{\fa{وثیقه} (vasiqe)}{وثیقه (vasiqe)}]{\fa{وثیقه} (vasiqe) — bail, a security}

\label{lesson:FA-C322-vasiqe}



\begin{quote}
Before the new one: say the Persian for to be acquitted, then the Persian for a prisoner.
\end{quote}

\subsection*{You'll want to know: \fa{وثیقه}}
\textbf{\fa{وثیقه}} (\emph{vasiqe}) — \textquotedblleft{}bail, a security\textquotedblright{}.

\begin{quote}
\fa{او} \fa{با} \fa{وثیقه} \fa{آزاد} \fa{شد}.
\end{quote}

\begin{quote}
\emph{U bâ vasiqe âzâd shod.}
\end{quote}

\begin{quote}
He was released on bail.
\end{quote}

Also a guarantee for a loan. Written with \fa{ث}.

\begin{cousinweb}
\emph{vasiqe} is usually traced to Arabic \emph{wathīqa}, \textquotedblleft{}a document, a pledge\textquotedblright{}; EtymDB lists no entry.
\end{cousinweb}

\begin{grammarlens}[title={What you've built}]
One more word for telling, explaining and describing at B1.
\end{grammarlens}

\subsection*{Your turn}
\begin{itemize}
\raggedright
  \item \emph{Say it:} \emph{vasiqe}
  \item \emph{Say it:} \emph{vasiqe}, once more
  \item \emph{Say it:} a sentence of your own with \emph{vasiqe}
  \item \emph{Recall:} say the Persian for to be acquitted, then the Persian for a prisoner, then say \emph{vasiqe} again
\end{itemize}

\begin{tcolorbox}[breakable,colback=teal!4,colframe=teal!35!black,title={Before you move on}]
What does \fa{وثیقه} mean? (\textquotedblleft{}Bail, a security\textquotedblright{}.) Say it once more.
\end{tcolorbox}
\section[\texorpdfstring{\fa{آزاد} \fa{شدن} (âzâd shodan)}{آزاد شدن (âzâd shodan)}]{\fa{آزاد} \fa{شدن} (âzâd shodan) — to be released, to be freed}

\label{lesson:FA-C322-azad-shodan}



\begin{quote}
Before the new one: say the Persian for a prisoner, then the Persian for bail.
\end{quote}

\subsection*{You'll want to know: \fa{آزاد} \fa{شدن}}
\textbf{\fa{آزاد} \fa{شدن}} (\emph{âzâd shodan}) — \textquotedblleft{}to be released, to be freed\textquotedblright{}.

\begin{quote}
\fa{زندانی} \fa{ها} \fa{صبح} \fa{آزاد} \fa{شدند}.
\end{quote}

\begin{quote}
\emph{Zendâni-hâ sobh âzâd shodand.}
\end{quote}

\begin{quote}
The prisoners were released in the morning.
\end{quote}

Present \emph{âzâd mishavad}. \emph{Âzâd kardan} is \textquotedblleft{}to set free\textquotedblright{}.

\begin{cousinweb}
\emph{âzâd} \textquotedblleft{}free\textquotedblright{} is inherited: EtymDB traces it through Middle Persian to Old Persian \emph{āzāta-}, \textquotedblleft{}noble, free\textquotedblright{}. \emph{shodan} \textquotedblleft{}to become\textquotedblright{} is inherited too.
\end{cousinweb}

\begin{grammarlens}[title={What you've built}]
One more word for telling, explaining and describing at B1.
\end{grammarlens}

\subsection*{Your turn}
\begin{itemize}
\raggedright
  \item \emph{Say it:} \emph{âzâd shodan}
  \item \emph{Say it:} \emph{âzâd shodan}, once more
  \item \emph{Say it:} a sentence of your own with \emph{âzâd shodan}
  \item \emph{Recall:} say the Persian for a prisoner, then the Persian for bail, then say \emph{âzâd shodan} again
\end{itemize}

\begin{tcolorbox}[breakable,colback=teal!4,colframe=teal!35!black,title={Before you move on}]
What does \fa{آزاد} \fa{شدن} mean? (\textquotedblleft{}To be released, to be freed\textquotedblright{}.) Say it once more.
\end{tcolorbox}
